\section{Hình học}
\begingroup\small
\begin{multicols}{2}

\subsection{Điểm, vector, tích vô hướng và tích có hướng}
Điểm $P=(x,y)$ cũng biểu diễn vector từ gốc tọa độ tới $P$.
Với $A,B$ là hai điểm, $\overrightarrow{AB}=B-A$.
Đường thẳng qua hai điểm \emph{phân biệt} $A,B$ có dạng tham số
$P(t)=A+t(B-A)$, $t\in\mathbb R$; đoạn $AB$ tương ứng
$0\leq t\leq1$. Đường tròn tâm $O$, bán kính $r\geq0$ gồm các
điểm thỏa $(x-O_x)^2+(y-O_y)^2=r^2$.

Với $u=(u_x,u_y)$ và $v=(v_x,v_y)$:
\[
u\cdot v=u_xv_x+u_yv_y,\qquad
u\times v=u_xv_y-u_yv_x.
\]
$u\cdot v$ cho góc, hình chiếu và khoảng cách;
$u\times v$ cho hướng quay và hai lần diện tích tam giác.
Đặt $\operatorname{orient}(A,B,C)=(B-A)\times(C-A)$:
giá trị dương là rẽ trái (ngược chiều kim đồng hồ), âm là rẽ phải,
bằng $0$ là thẳng hàng. Với tọa độ nguyên, hãy xét dấu \emph{chính xác}
thay vì dùng $\varepsilon$.

\subsection{Đường thẳng, đoạn thẳng và giao điểm}
Đường thẳng qua $P_1=(x_1,y_1)$, $P_2=(x_2,y_2)$ cũng có thể viết
$Ax+By=C$, trong đó
\[
A=y_2-y_1,\quad B=x_1-x_2,\quad C=Ax_1+By_1.
\]
Với $L_1=P+tu$, $L_2=Q+sv$, nếu $u\times v\ne0$ thì giao điểm duy
nhất là
\[
P+u\,\frac{(Q-P)\times v}{u\times v}.
\]
Nếu $u\times v=0$, hai đường song song hoặc trùng nhau; kiểm tra thêm
$(Q-P)\times u=0$ để nhận biết trùng nhau. Không áp dụng công thức
cho đường có vector chỉ phương bằng $0$.

Đường thẳng $AB$ cắt đoạn $CD$ tại một điểm nếu $C,D$ ở hai phía
khác nhau hoặc có đầu mút trên đường; nếu cả hai cùng nằm trên
đường thì giao là cả đoạn (hoặc một điểm nếu $C=D$). Với
\emph{hai đoạn}, trước hết xét dấu
$\operatorname{orient}$ của mỗi cặp đầu mút đối với đường chứa
đoạn kia; trường hợp bằng $0$ phải kiểm tra đầu mút có nằm trong
hộp tọa độ của đoạn hay không. Code bên dưới tính cả trường hợp
chạm đầu mút, chồng đoạn và đoạn chỉ có một điểm.

\subsection{Khoảng cách}
Với $u=B-A\ne0$, khoảng cách từ $Q$ tới đường thẳng $AB$ là
\[
d(Q,AB)=\frac{|u\times(Q-A)|}{\lVert u\rVert}.
\]
Với \emph{đoạn} $AB$, lấy $t=((Q-A)\cdot u)/(u\cdot u)$ rồi chặn
$t$ vào $[0,1]$; điểm gần nhất là $A+tu$. Nếu $A=B$, khoảng cách
là $\lVert Q-A\rVert$. So sánh khoảng cách bình phương khi không
cần chính giá trị căn bậc hai.

\subsection{Diện tích đa giác đơn}
Cho đa giác \emph{không tự cắt} có $n$ đỉnh $P_0,\ldots,P_{n-1}$
theo thứ tự quanh biên; đặt $P_n=P_0$. Hai cách viết tương đương của
\emph{hai lần diện tích có dấu} là
\[
S_2=\sum_{i=0}^{n-1}(x_i y_{i+1}-x_{i+1}y_i)
=\sum_{i=0}^{n-1}(x_i-x_{i+1})(y_i+y_{i+1}).
\]
Hai tổng bằng nhau vì các hạng $x_i y_i$ triệt tiêu khi lấy tổng
theo chu trình. $S_2>0$ khi duyệt ngược chiều kim đồng hồ, $S_2<0$
khi duyệt thuận chiều; diện tích là $|S_2|/2$. Với tọa độ nguyên,
$S_2$ nguyên nên diện tích là số nguyên hoặc số nguyên cộng $0.5$.
Nếu đa giác tự cắt, $|S_2|/2$ nói chung không phải diện tích miền.

\subsection{Bao lồi: Graham scan}
Bao lồi là tập lồi nhỏ nhất chứa mọi điểm; nếu không suy biến,
nó là một đa giác lồi. Graham scan: bỏ điểm
trùng, chọn đỉnh thấp nhất (nếu hòa chọn trái nhất), sắp xếp các
điểm còn lại theo góc cực quanh đỉnh đó, rồi duy trì một stack;
khi ba điểm cuối không rẽ trái thì bỏ điểm giữa. Với các điểm cùng
tia, sắp xếp gần trước và loại điểm gần. Độ phức tạp $O(n\log n)$;
code trả các đỉnh theo chiều ngược kim đồng hồ, không lặp đỉnh đầu,
và bỏ các điểm thẳng hàng nằm giữa trên cạnh. Nếu mọi điểm thẳng
hàng, kết quả là hai đầu mút (hoặc một điểm nếu tất cả trùng nhau).

\paragraph{Cài đặt số nguyên (C++17).}
Các phép nhân và hiệu tọa độ được nâng lên \texttt{\_\_int128}
trước khi tính; vẫn cần kiểm tra giới hạn dữ liệu để tổng cuối không
vượt kiểu này. \texttt{polygonArea2} dùng đúng công thức $S_2$ ở trên.
\cppsnippet{code/geometry/integer.hpp}

\subsection{Đường tròn}
Để cắt đường thẳng $L=P+tu$ với đường tròn $(O,r)$, lấy hình chiếu
$H=P+u\,((O-P)\cdot u)/(u\cdot u)$ và $d^2=\lVert H-O\rVert^2$.
Nếu $d^2>r^2$ không có giao; bằng nhau thì tiếp xúc tại $H$;
nhỏ hơn thì hai giao điểm là
\[
H\pm u\sqrt{\frac{r^2-d^2}{u\cdot u}}.
\]
Muốn giao với \emph{đoạn}, chỉ giữ các nghiệm có tham số $t\in[0,1]$.

\paragraph{Cài đặt số thực (C++17).}
\texttt{Line} dùng điểm $p$ và vector hướng $v$; \texttt{Circle} dùng
tâm và bán kính. Hàm giao hai đường trả rỗng cho cả song song lẫn
trùng nhau, nên cần kiểm tra riêng nếu bài đòi phân biệt hai trường
hợp. Vector hướng của đường phải khác $0$.
\cppsnippet{code/geometry/real.hpp}

\subsection{Sai số số thực và các trường hợp suy biến}
Dùng phép tính nguyên và \texttt{\_\_int128} cho hướng quay, giao
đoạn, bao lồi và diện tích khi tọa độ nguyên; kiểm tra giới hạn
để tránh tràn ngay cả với kiểu này. Khi cần giao điểm hoặc căn bậc
hai, dùng \texttt{long double} và chọn $\varepsilon$ theo thang tọa độ
của bài. Không dùng so sánh có $\varepsilon$ trong comparator của
\texttt{sort}, vì có thể phá tính bắc cầu. Cần xét riêng đường thẳng
không hợp lệ, đoạn độ dài $0$, điểm thẳng hàng và trường hợp tiếp xúc.

\end{multicols}
\endgroup